\documentclass[../../main.tex]{subfiles}

\begin{document}

\section{自反 Banach 空间与弱序列收敛性}\label{section:ch14-自反弱序列}

\begin{theorem}\label{theorem:ch14-对偶可分}
令 $X$ 为赋范线性空间. 若它的对偶空间 $X^{*}$ 是可分的, 则 $X$ 也是可分的.
\end{theorem}

\begin{proof}
由于 $X^{*}$ 是可分的, 它的闭单位球 $S^{*}=\{\psi\in X^{*}\mid\|\psi\|=1\}$ 也如此. 令 $\{\psi_{n}\}_{n=1}^{\infty}$ 为 $S^{*}$ 的可数稠密子集. 对每个指标 $n$, 选取 $x_{n}\in X$ 使得:
\begin{align*}
\|x_{n}\|=1\text{ 且 }\psi_{n}(x_{n})>1/2.
\end{align*}
定义 $X_{0}$ 为集合 $\{x_{n}\mid 1\leqslant n<\infty\}$ 的闭线性扩张. 则 $X_{0}$ 是可分的, 由于具有有理系数的 $x_{n}$ 的有限线性组合是 $X_{0}$ 的可数稠密子集. 我们宣称 $X_{0}=X$. 事实上, 若不是这样, 根据定理 10, 我们可以选取 $\psi^{*}\in X^{*}$ 使得:
\begin{align*}
\|\psi^{*}\|=1\text{ 且在 }X_{0}\text{ 上 }\psi^{*}=0.
\end{align*}
由于 $\{\psi_{n}\}_{n=1}^{\infty}$ 在 $S^{*}$ 中稠密, 存在自然数 $n_{0}$, 使得 $\|\psi^{*}-\psi_{n_{0}}\|<1/2$. 因此
\begin{align*}
|(\psi_{n_{0}}-\psi^{*})(x_{n_{0}})|\leqslant\|\psi_{n_{0}}-\psi^{*}\|\cdot\|x_{n_{0}}\|<1/2.
\end{align*}
然而有 $(\psi_{n_{0}}-\psi^{*})(x_{n_{0}})=\psi_{n_{0}}(x_{n_{0}})>1/2$. 从这个矛盾我们推出 $X$ 是可分的.
\end{proof}

\begin{corollary}\label{corollary:ch14-自反可分}
自反 Banach 空间是可分的当且仅当它的对偶是可分的.
\end{corollary}

\begin{proof}
令 $X$ 为 Banach 空间. 前一个定理告诉我们, 若 $X^{*}$ 是可分的, 则 $X$ 也是可分的, 与任何自反性假设无关. 现在假设 $X$ 是自反的与可分的. 因此 $J(X)=X^{**}=(X^{*})^{*}$ 是可分的, 由于 $J$ 是等距. 根据前一个定理, 其中 $X$ 用 $X^{*}$ 代替, $X^{*}$ 是可分的.
\end{proof}

\begin{proposition}\label{proposition:ch14-自反子空间}
自反 Banach 空间的闭子空间是自反的.
\end{proposition}

\begin{proof}
令 $X_{0}$ 为自反 Banach 空间 $X$ 的闭子空间. 定义 $J$ 为 $X$ 到它的二次对偶 $X^{**}$ 的自然嵌入. $J_{0}$ 为 $X_{0}$ 到它的二次对偶在 $X_{0}^{**}$ 的自然嵌入. 为证明 $J_{0}$ 是映上的, 令 $S$ 属于 $X_{0}^{**}$. 定义 $S'\in X^{**}$ 为:
\begin{align*}
\text{对所有 }\psi\in X^{*},\quad S'(\psi)=S(\psi|_{X_{0}}).
\end{align*}
则 $S':X_{0}^{*}\to\mathbf{R}$ 是线性的且它是有界的, 满足 $\|S'\|\leqslant\|S\|$. 根据 $X$ 的自反性, 存在 $x_{0}\in X$ 使得 $S'=J(x_{0})$. 但若 $\psi\in X^{*}$ 在 $X_{0}$ 上消失, 则 $S'(\psi)=0$, 因此,
\begin{align*}
\psi(x_{0})=J(x_{0})[\psi]=S'(\psi)=0.
\end{align*}
定理 10 告诉我们, $x_{0}$ 属于 $X_{0}$. 因此 $S=J_{0}(x_{0})$.
\end{proof}

我们再次写出在第 8 章证明过的 Helly 定理.

\begin{theorem}[Helly 定理]\label{theorem:ch14-Helly}
令 $X$ 为可分赋范线性空间. 则 $X^{*}$ 中的每个有界序列 $\{\psi_{n}\}$ 有在 $X$ 上逐点收敛于 $\psi\in X^{*}$ 的子序列, 即 $\{\psi_{n}\}$ 有关于弱$^{*}$拓扑收敛到 $\psi$ 的子序列.
\end{theorem}

\begin{theorem}\label{theorem:ch14-自反弱序列紧}
令 $X$ 为自反 Banach 空间. 则 $X$ 中的每个有界序列有弱收敛的子序列.
\end{theorem}

\begin{proof}
令 $\{x_{n}\}$ 为 $X$ 中的有界序列. 定义 $X_{0}$ 为集合 $\{x_{n}\mid n\in\mathbf{N}\}$ 的线性扩张的闭包. 则 $X_{0}$ 是可分的, 由于 $x_{n}$ 的具有有理系数的有限线性组合是 $X_{0}$ 的可数稠密子集. 当然 $X_{0}$ 是闭的. 命题 15 告诉我们 $X_{0}$ 是自反的. 令 $J_{0}$ 为 $X_{0}$ 到它的二次对偶 $X_{0}^{**}$ 的自然嵌入. 从推论 14 得出 $X_{0}^{*}$ 也是可分的. 则 $\{J_{0}(x_{n})\}$ 是可分 Banach 空间 $X_{0}^{*}$ 上的有界线性泛函的有界序列. 根据 Helly 定理, 子序列 $\{J_{0}(x_{n_{k}})\}$ 弱$^{*}$收敛于 $S\in(X_{0}^{*})^{*}$. 由于 $X_{0}$ 是自反的, 存在某个 $x_{0}\in X_{0}$ 使得 $S=J_{0}(x_{0})$. 由于 $X^{*}$ 中的每个泛函可限制为 $X_{0}^{*}$ 上的泛函, $\{J_{0}(x_{n_{k}})\}$ 弱$^{*}$收敛于 $J_{0}(x_{0})$ 恰好意味着 $\{x_{n_{k}}\}$ 弱收敛于 $x_{0}$.
\end{proof}

\begin{corollary}\label{corollary:ch14-自反取最大}
令 $X$ 为自反 Banach 空间. 则 $X$ 上的每个连续实值线性泛函在 $X$ 的闭单位球 $B$ 上取到最大值.
\end{corollary}

\begin{proof}
令 $\psi$ 属于 $X^{*}$. $\psi$ 在 $B$ 上的泛函值的上确界是 $\|\psi\|$. 选取 $B$ 中的序列 $\{x_{n}\}$ 使得 $\lim_{n\to\infty}\psi(x_{n})=\|\psi\|$. 根据定理 17, 我们可以假设 $\{x_{n}\}$ 弱收敛到 $x_{0}$. 根据\eqref{align:ch14-弱收敛范数}, $x_{0}$ 属于 $B$. 由于
\begin{align*}
\lim_{n\to\infty}\psi(x_{n})=\psi(x_{0}),
\end{align*}
$\psi$ 在 $x_{0}$ 取到 $B$ 上的最大值.
\end{proof}

定理 17 使得识别哪些经典 Banach 空间是自反的是有趣的.

\begin{proposition}\label{proposition:ch14-C[a,b]不自反}
令 $[a,b]$ 为实数的闭有界区间. 则赋予最大值范数的 $C[a,b]$ 不是自反的.
\end{proposition}

\begin{proof}
假设 $[a,b]=[0,1]$. 对 $x\in[0,1]$, 定义估值函数 $\psi_{x}:C[0,1]\to\mathbf{R}$ 为 $\psi_{x}(f)=f(x)$. 则 $\psi_{x}$ 是 $C[0,1]$ 上的有界线性泛函. 因此, 若 $\{f_{n}\}$ 在 $C[0,1]$ 上弱收敛于 $f$, 则在 $[0,1]$ 上逐点地 $\{f_{n}\}\to f$. 对自然数 $n$, 定义 $f_{n}(x)=x^{n}$, $x\in[0,1]$. 则 $\{f_{n}\}$ 逐点收敛到不连续函数 $f$. 因此没有子序列逐点收敛到连续函数, 且因此没有子序列弱收敛到 $C[0,1]$ 中的函数. 我们从定理 17 推出, $C[0,1]$ 不是自反的.
\end{proof}

\begin{proposition}\label{proposition:ch14-Lp自反}
对 $1<p<\infty$ 而 $E$ 是实数的 Lebesgue 可测集, $L^{p}(E)$ 是自反的.
\end{proposition}

\begin{proof}
令 $q$ 为 $p$ 的共轭. 定义从 $L^{q}(E)$ 到 $(L^{p}(E))^{*}$ 的 Riesz 表示映射 $\mathcal{R}$ 为
\begin{align*}
\text{对所有 }g\in L^{q}(E),\ f\in L^{p}(E),\quad\mathcal{R}(g)[f]=\int_{E}g\cdot f.
\end{align*}
Riesz 表示定理告诉我们 $\mathcal{R}$ 是 $L^{q}(E)$ 映上 $(L^{p}(E))^{*}$ 的同构. 令 $T$ 为 $(L^{p}(E))^{*}$ 上的有界线性泛函. 我们必须证明存在函数 $f\in L^{p}(E)$ 使得 $T=J(f)$, 即由于 $\mathcal{R}$ 是映上的,
\begin{align}
\text{对所有 }g\in L^{q}(E),\quad T(\mathcal{R}(g))=J(f)[\mathcal{R}(g)]=\mathcal{R}(g)[f]=\int_{E}g\cdot f.\label{align:ch14-Lp自反}
\end{align}
然而, 复合 $T\circ\mathcal{R}$ 是 $L^{q}(E)$ 上的有界线性泛函. Riesz 表示定理(其中 $p,q$ 互换)告诉我们, 存在函数 $f\in L^{p}(E)$ 使得:
\begin{align*}
\text{对所有 }g\in L^{q}(E),\quad(T\circ\mathcal{R})[g]=\int_{E}f\cdot g.
\end{align*}
因此,
\begin{align*}
\text{对所有 }g\in L^{q}(E),\quad T(\mathcal{R}(g))=(T\circ\mathcal{R})[g]=\int_{E}f\cdot g=\int_{E}g\cdot f.
\end{align*}
即\eqref{align:ch14-Lp自反}成立.
\end{proof}

一般来说, $L^{1}(E)$ 与 $L^{\infty}(E)$ 都不是自反的. 考虑 $E=[0,1]$. 观察到 $L^{1}[0,1]$ 是可分的而 $(L^{1}[0,1])^{*}$ 不是可分的, 由于它同构于不可分的 $L^{\infty}[0,1]$. 我们从推论 14 推出 $L^{1}[0,1]$ 不是自反的. 观察到 $C[0,1]$ 是 $L^{\infty}[0,1]$ 的闭子空间(见习题 29). 根据命题 19, $C[0,1]$ 不是自反的, 从而命题 15 告诉我们, $L^{\infty}[0,1]$ 也不是自反的.

\textbf{注} 当证明自反性时需要注意. R. C. James 已给出同构于 $X^{**}$ 但不自反的 Banach 空间 $X$ 的例子. 自反性不仅要求 $X$ 同构于 $X^{**}$, 它要求 $X$ 到 $X^{**}$ 的自然嵌入 $J$ 是一个同构映射.

\textbf{注} 可积函数空间与连续函数空间的自反性之间的反差是惊人的. 第 19 章的定理 8 告诉我们, 对 $1<p<\infty$, 最一般的 $L^{p}$ 空间是自反的. 另一方面, 若 $K$ 是任何紧 Hausdorff 空间而 $C(K)$ 赋予最大值范数, 则 $C(K)$ 是自反的当且仅当 $K$ 是有限集(见第 15 章的习题 11).

\begin{exercise}\label{exercise:ch14-31}
证明有界线性函数族是等度连续的当且仅当它是一致有界的.
\end{exercise}

\begin{exercise}\label{exercise:ch14-32}
令 $X$ 为可分赋范线性空间. 证明它的闭单位球 $S=\{x\in X\mid\|x\|=1\}$ 也是可分的.
\end{exercise}

\begin{exercise}\label{exercise:ch14-33}
找到紧度量空间 $X$, 使得赋予最大值范数的 $C(X)$ 是自反的.
\end{exercise}

\begin{exercise}\label{exercise:ch14-34}
令 $c_{0}$ 为由收敛到 $0$ 的序列组成的 $\ell^{\infty}$ 的子空间. 证明 $c_{0}$ 是 $\ell^{\infty}$ 的闭子空间, 其对偶空间同构于 $\ell^{1}$. 得出 $c_{0}$ 不是自反的, 因此 $\ell^{\infty}$ 也不是自反的.
\end{exercise}

\begin{exercise}\label{exercise:ch14-35}
对 $1\leqslant p\leqslant\infty$, 证明序列空间 $\ell^{p}$ 是自反的当且仅当 $1<p<\infty$. (关于 $p=\infty$, 见前一个习题.)
\end{exercise}

\begin{exercise}\label{exercise:ch14-36}
考虑定义为
\begin{align*}
\text{对所有 }h\in C[-1,1],\quad\psi(h)=\int_{-1}^{0}h-\int_{0}^{1}h
\end{align*}
的泛函 $\psi\in(C[-1,1])^{*}$. 证明 $\psi$ 在 $C[-1,1]$ 的闭单位球上取不到最大值. 用此给出 $C[-1,1]$ 不是自反的另一证明.
\end{exercise}

\begin{exercise}\label{exercise:ch14-37}
对 $1<p<\infty$, 证明 $\ell^{p}$ 中的有界序列弱收敛当且仅当它依分量收敛.
\end{exercise}

\begin{exercise}\label{exercise:ch14-38}
对 $1\leqslant p\leqslant\infty$, $[a,b]$ 是实数的闭有界区间, 证明 $L^{p}[a,b]$ 中的有界序列 $\{f_{n}\}$ 弱收敛到 $f$, 当且仅当对 $[a,b]$ 的每个 Lebesgue 可测子集 $E$,
\begin{align*}
\left\{\int_{E}f_{n}\right\}\to\int_{E}f.
\end{align*}
\end{exercise}

\begin{exercise}\label{exercise:ch14-39}
对实数的闭有界区间 $[a,b]$, 证明: 若 $C[a,b]$ 中的序列 $\{f_{n}\}$ 弱收敛, 则它逐点收敛.
\end{exercise}

\begin{exercise}\label{exercise:ch14-40}
$X$ 和 $Y$ 为赋范线性空间而算子 $S\in\mathcal{L}(X,Y)$, 定义 $S$ 的伴随 $S^{*}\in\mathcal{L}(Y^{*},X^{*})$ 为
\begin{align*}
\text{对所有 }\psi\in Y^{*},\ x\in X,\quad[S^{*}(\psi)](x)=\psi(S(x)).
\end{align*}
(i)证明 $\|S^{*}\|=\|S\|$, 且若 $S$ 是同构则 $S^{*}$ 是同构.
(ii)对 $1<p<\infty$ 与 $X=L^{p}(E)$, 其中 $E$ 是实数的可测集, 证明自然嵌入 $J:X\to X^{**}$ 可表示为复合
\begin{align*}
J=[\mathcal{R}_{q}^{*}]^{-1}\circ\mathcal{R}_{p}
\end{align*}
其中 $\mathcal{R}_{p}$ 和 $\mathcal{R}_{q}$ 是 Riesz 表示算子.
\end{exercise}

\begin{exercise}\label{exercise:ch14-41}
令 $X$ 为自反 Banach 空间而 $T:X\to X$ 是线性算子. 证明 $T$ 属于 $\mathcal{L}(X,X)$ 当且仅当若 $\{x_{n}\}$ 弱收敛于 $x$, 则 $\{T(x_{n})\}$ 弱收敛于 $T(x)$.
\end{exercise}

\end{document}